\chapter{实验验证与系统评估}\label{chap5_experiment_evaluation}

\section{实验平台与典型工况设计}\label{sec_ch5_platform_scene}

\subsection{吊钩端实验平台与传感器组成}\label{subsec_ch5_platform}

第4章给出了从 YOLO26 分割识别到点云投影定位、再到预警与避障设计的完整技术路线。为了验证该路线在真实吊装场景中的可行性，本章首先说明实验平台的组成及其运行方式。当前实验平台以吊钩端一体化感知盒体为核心，集成工业相机、Livox 激光雷达、IMU、Jetson 边缘计算单元及 ROS 运行环境。视觉侧负责吊装构件图像采集与 YOLO26 分割推理，LiDAR 侧负责提供三维点云及其动态信息，定位与跟踪模块则负责坐标统一、动态候选提取和轨迹预测。

从系统组织形式看，整个平台并不是多个独立算法的简单叠加，而是围绕“感知输入统一、状态输出统一、预警链路可承接”的原则构建。视觉分割节点输出吊装构件的实例掩码与可视化结果，点云投影模块输出语义着色点云与三维质心，动态检测链路输出动态障碍物候选与预测轨迹，最终这些结果在统一坐标参考下进入预警模块。这样的平台组织方式保证了实验章节不是孤立评价某一个感知模块，而是围绕整条感知到预警链路展开。

对于论文写作而言，本小节还需要强调平台构成与章节评价指标之间的对应关系。视觉识别效果评价依赖相机和 YOLO26 推理链路，空间定位一致性评价依赖点云投影与三维质心提取，动态预警有效性评价依赖动态点检测与卡尔曼预测，端到端实时性分析则依赖系统各阶段的耗时记录。因此，本章实验设计从一开始就不是“把所有模块都测一遍”，而是要验证第4章提出的方法链路在平台上能否形成稳定、可解释、可扩展的工程能力。

\subsection{典型测试场景与rosbag数据组织}\label{subsec_ch5_scenes_bag}

为使实验结果具有代表性，本研究围绕吊装作业中的风险形成机制设计典型工况。总体上，测试场景至少应覆盖四类情况：第一类是吊装构件与静态障碍物逐步接近的场景，用于验证空间距离阈值预警的有效性；第二类是动态人员或车辆在吊物附近运动的场景，用于验证动态候选生成、轨迹预测和风险升级判据的合理性；第三类是吊钩或吊物自身存在摆动和姿态变化的场景，用于验证图像与点云同步、投影稳定性以及三维定位的鲁棒性；第四类是吊物与动态目标相向接近的场景，用于重点检验预警提前量与分级动作策略。

从实验组织方式看，在线采集与离线回放应当结合使用。在线实验适合观察系统在真实运行状态下的模块协同表现，离线 rosbag 回放则更适合在同一数据条件下反复调整参数、复现实验过程和统计评价指标。对于本研究而言，已有多种典型场景数据可为静态障碍、动态目标接近和吊钩摆动等实验提供复现基础，但其具体路径和场景标注尚需进一步整理，因此论文中可写为 \textbf{【待补充：rosbag 数据集清单、时间长度与场景标签】}。

本小节的重点并不只是罗列若干测试名称，而是说明实验如何覆盖第4章各个关键模块的风险边界。静态场景主要服务于空间定位与距离阈值卡控验证，动态场景主要服务于动态趋势预测和预警升级验证，摆动场景则检验投影与同步链路在复杂工况下的稳定性。这样安排的结果是，第5章的实验内容能够与第4章的方法结构一一对应，避免出现“方法写的是吊装构件定位与预警，实验测的却仍是旧版车辆检测流程”的断裂问题。

\begin{figure}[htbp]
	\centering
	\fbox{\parbox[c][5.4cm][c]{0.9\textwidth}{\centering
	占位图：实验平台与典型工况组织示意图\\[0.5em]
	建议内容：吊钩端平台、相机、双 LiDAR、Jetson、ROS 链路，以及静态障碍物场景、动态接近场景、吊物摆动场景和相向运动场景的示意分图。}}
	\caption{实验平台与典型工况组织占位图}
	\label{fig:ch5_scene_design_placeholder}
\end{figure}

\section{吊装构件识别与空间定位效果分析}\label{sec_ch5_recognition_localization}

\subsection{YOLO26吊装构件检测与分割效果评估}\label{subsec_ch5_yolo26_eval}

视觉侧实验首先关注 YOLO26 在吊装场景中的识别与分割表现。与旧版以车辆为代表目标的检测口径不同，本章评价对象明确调整为吊装构件或吊物本身，即模型是否能够在复杂工地背景下稳定给出目标区域的分割结果。该评价的意义在于，第4章后续所有点云投影与空间定位步骤都依赖分割掩码的空间约束能力，若分割边界质量不足，则后续点云回查会受到直接影响。

在评价内容上，建议从三个层面展开。第一是\textbf{识别有效性}，即模型能否在工地实采数据中稳定检测到吊装构件目标，可用召回率、漏检率和持续漏检帧数等指标衡量。第二是\textbf{分割边界质量}，即预测掩码与真实目标区域在边界上的一致程度，可进一步采用 IoU、Dice 系数或人工核查结果进行分析。第三是\textbf{复杂背景鲁棒性}，即在遮挡、强光、阴影、构件尺度变化和背景杂乱条件下，模型输出是否仍保持可用。由于当前仓库中尚未整理完整训练验证结果，因此论文正文可先把评价口径写清，并将具体数值留作 \textbf{【待补充：YOLO26 分割精度与场景分组结果】}。

需要强调的是，本小节的实验目标不是证明 YOLO26 在所有通用视觉任务中都优于其他模型，而是验证其是否满足本文工程任务要求，即能否在边缘端稳定提供可用于点云回查的吊装构件区域。只要模型在主要场景下能够稳定输出合理掩码，并在时延上满足在线要求，它就已经具备支撑后续空间定位与预警链路的条件。因此，本小节的结论应更多地围绕“是否足够支撑系统任务”展开，而不是单纯追求模型排行榜意义上的精度比较。

\subsection{点云投影关联与三维定位一致性验证}\label{subsec_ch5_projection_localization_eval}

在视觉识别结果具备可用性的基础上，实验还需要进一步验证分割掩码与点云投影之间的空间对应关系是否稳定，以及吊装构件的三维定位结果是否具有一致性。该部分评价直接对应第4.2节的核心方法，因为点云投影回查和三维质心提取是把二维语义结果转化为三维安全基准的关键步骤。若这一环节不稳定，则后续预警阈值计算将缺少可靠的空间基础。

本小节建议围绕三个问题展开。第一，\textbf{投影重合度是否稳定}，即投影到图像平面的 LiDAR 点是否能够持续落入正确的分割掩码区域，特别是在吊物摆动、姿态变化和局部遮挡条件下，投影结果是否仍与目标轮廓保持一致。第二，\textbf{三维质心是否平滑可信}，即由实例点集得到的质心位置在连续帧中是否具有合理变化趋势，是否会因少量噪声点或掩码边界抖动而发生明显跳变。第三，\textbf{不同姿态和距离下的空间定位误差}，即当吊装构件处于不同视角、不同距离和不同遮挡程度时，三维定位结果与人工测量或先验标定位置之间的偏差大小。

目前系统已经能够输出投影叠加图像、语义着色点云和三维质心信息，因此该实验具有较强的可操作性。对于尚未补齐的定量误差结果，可以在正文中明确写成 \textbf{【待补充：投影重合度统计、质心抖动范围与定位误差】}。这样既保留了第4章方法已经具备验证基础这一事实，也为后续补充更完整的实验曲线、表格和案例图预留空间。本小节最终要回答的问题不是“视觉和点云是否能勉强对上”，而是“这种对齐是否稳定到足以作为预警距离计算的基准”。

\section{预警链路有效性实验}\label{sec_ch5_warning_validation}

\subsection{静态障碍物距离阈值预警效果测试}\label{subsec_ch5_static_warning}

静态障碍物预警实验对应第4.3.1节提出的分层距离阈值卡控方法，目的是验证吊装构件空间定位结果能否稳定支撑多级安全边界判断。实验中应在吊物运动路径附近布置典型静态障碍物，如支撑结构、脚手架边缘或固定构件，并通过改变目标与障碍物之间的最近距离，观察系统输出的距离估计与预警等级变化。由于此处的核心问题是“距离阈值是否可靠”，因此实验应尽量避免动态因素干扰，使结果能够直接反映空间定位与阈值设计之间的关系。

评价指标建议包括四项。第一是\textbf{触发距离误差}，即系统给出的预警触发时刻对应的距离与人工测量值之间的差异。第二是\textbf{分级稳定性}，即一级、二级和三级预警是否按设定阈值依次触发，是否存在频繁跨级跳变。第三是\textbf{误报与漏报}，即在明显安全距离下是否仍触发预警，以及在明显危险接近时是否出现未触发现象。第四是\textbf{响应延迟}，即从空间关系实际进入某一阈值区域到系统发布相应等级预警之间的时间差。上述指标可以较好地反映静态阈值卡控作为基础预警层的有效性。

需要说明的是，当前仓库中尚未形成完整的静态分层预警模块，因此本实验在论文中应写成“基于新预警模块或离线分析脚本的验证方案”。也就是说，第4章给出了明确的风险卡控设计，第5章则给出相应的实验验证口径和结果组织方式；具体数值统计可在后续补充为 \textbf{【待补充：各等级阈值、触发误差和误报漏报统计】}。这种写法能够保持方法与实验的一致性，同时避免把设计性内容误写成已经完全落地的现成功能。

\subsection{动态目标趋势预测与相向运动预警验证}\label{subsec_ch5_dynamic_warning}

动态场景实验的重点在于验证卡尔曼预测和风险升级规则是否能够提供比单一距离阈值更合理的提前量。具体而言，应构造动态目标与吊装构件存在横向通过、斜向接近和相向运动等不同关系的场景，比较系统在这些工况下给出的预测轨迹、最近接距离估计和预警等级变化。若仅采用瞬时距离阈值，则系统在相向运动场景中的预警通常偏晚；而引入趋势预测之后，理论上可以在距离尚未极度危险时提前进入较高等级，从而为控制动作留出反应时间。

实验评价可围绕三类结果展开。第一类是\textbf{轨迹预测一致性}，即卡尔曼预测轨迹是否与目标真实运动趋势保持一致，可通过轨迹可视化对比和短时预测误差统计进行分析。第二类是\textbf{风险升级合理性}，即相向运动场景是否能够比横向远离场景更早触发高等级预警，且不会因短时噪声导致频繁误升。第三类是\textbf{预警提前量}，即系统在动态接近过程中相较于纯距离阈值方案能够提前多少时间或多少距离发出报警。该部分是验证第4.3.4节设计价值的关键实验。

由于当前系统已具备动态点检测与卡尔曼预测基础，但 TTC 计算和完整三级状态机仍属于方法设计，因此本小节可采用“现有预测模块 + 离线风险评估”的方式组织实验。正文中可写明“动态候选生成与趋势预测已有实现基础，风险升级规则的量化结果记为 \textbf{【待补充：预测误差、相向场景提前量与等级切换统计】}”。这种安排能够准确呈现当前研究状态，即预警升级逻辑已具备可验证前提，但仍需通过更完整的实验统计加以支撑。

\subsection{预警到控制输出的动作响应验证}\label{subsec_ch5_control_response}

为了验证第4.4节提出的控制接口与动作策略是否具备原型可行性，本章还需要设计从预警等级到控制建议输出的动作响应实验。该实验并不要求系统已经接入真实塔吊控制器，而是要验证在给定风险等级条件下，控制接口是否能够稳定生成与设计一致的动作建议，例如一级提示对应观察，二级报警对应限速或限向，三级紧急对应停止建议。通过这种方式，可以在控制闭环尚未实机接入的阶段，先完成接口级和逻辑级的正确性验证。

实验可采用离线回放和原型仿真两种方式进行。离线回放适合验证不同预警场景下输出字段是否完整、动作类型是否匹配；原型仿真则适合进一步观察若控制建议被执行，系统是否能够在逻辑上形成闭环响应。评价指标可包括命令生成正确率、动作触发延迟、等级与动作映射一致性以及异常场景下的保护逻辑触发情况。由于当前仓库中尚未实现完整控制接口，因此本小节应明确表述为“设计验证或原型验证”，而不能写成已完成实机动作闭环。

从论文结构上看，本小节的作用是把第4章中的避障设计从纯概念进一步推进到“可验证接口”层面。即使当前实验仍停留在离线或原型阶段，它也已经表明本文研究并非只关注前端感知，而是把结果如何作用于后续安全辅助决策一并纳入了系统评估范围。具体数值可留作 \textbf{【待补充：控制输出正确率与响应时间】}。

\section{整体链路时延分析与实时性评估}\label{sec_ch5_latency}

\subsection{视觉推理模块时延分析}\label{subsec_ch5_infer_latency}

视觉推理时延是整条感知链路的首要实时性指标，因为 YOLO26 的前向推理直接决定了后续投影、定位和预警更新的起始节奏。当前系统已经在视觉节点中记录单帧推理耗时，因此本小节可以围绕 infer\_ms 的统计结果展开。分析时不应只关注平均值，还应关注最大值、方差和高分位数，因为吊装预警系统对时延抖动同样敏感。若极端帧的推理耗时过大，即使均值满足要求，系统仍可能在动态接近场景下出现预警迟滞。

从工程角度看，FP16 部署对视觉实时性的贡献应当在本小节中得到体现。可通过对比不同输入分辨率、不同场景复杂度或不同运行负载下的推理耗时变化，分析边缘端推理瓶颈是否主要来自网络计算本身，还是来自数据搬运与后处理。由于当前仓库尚未整理完整实验表格，因此论文可先建立明确的统计口径，并将具体数值写为 \textbf{【待补充：视觉推理均值、方差及 95\% 分位】}。这种处理既保留了当前计时功能已经存在的事实，也为后续填充正式实验数据提供直接框架。

\subsection{点云投影与图像关联模块时延分析}\label{subsec_ch5_projection_latency}

与单独的视觉推理相比，点云投影与图像关联模块更能反映多传感器融合系统的工程代价。当前视觉节点中已经对点云转换、下采样、TF 查询、投影绘制、叠加显示、结果发布以及总耗时进行了阶段计时，因此本小节具备较明确的量化基础。通过这些阶段统计，可以把“投影关联是否能在线运行”这一问题拆解为多个可观察的子问题，而不是只给出一个模糊的总时延数字。

本小节建议重点分析三部分。第一是\textbf{点云预处理开销}，包括格式转换和下采样，其结果直接影响单帧可处理点数和后续投影复杂度。第二是\textbf{变换与投影开销}，包括 TF 查询和 3D 到 2D 的投影计算，这部分反映了同步、坐标变换和像素级回查对系统带来的额外代价。第三是\textbf{结果组织与发布开销}，包括叠加绘制与消息发布，其结果决定了系统是否还能在可视化和调试需求下保持实时性。若后续需要进一步压缩总时延，本小节的阶段分析也能够为优化方向提供依据。

此外，近似时间同步窗口和下采样参数也应作为时延分析中的控制变量。同步窗口过大可能提高可匹配率，但也会放大错帧风险；下采样过强则可降低计算负担，但可能损失目标点集的空间密度。因此，本小节不仅要回答“模块耗时是多少”，还要回答“哪些参数变化会显著影响实时性与投影稳定性之间的平衡”。对应的实验统计可在后续补充为 \textbf{【待补充：各阶段耗时表与参数敏感性分析】}。

\subsection{动态跟踪与风险评估模块时延分析}\label{subsec_ch5_dynamic_latency}

除视觉投影链路外，动态障碍物检测、聚类、轨迹预测和风险评估也会共同占用系统时延预算。当前代码在动态跟踪方面已经形成完整处理链路，但尚未像视觉模块那样提供细粒度阶段计时，因此本小节更适合作为“基于时间戳统计的补充评估框架”来展开。其目的不是强行给出当前仓库没有的精确数值，而是说明在完整系统评估中，这一部分必须被单独量化。

在分析口径上，可将动态链路拆分为动态点提取、聚类中心计算、卡尔曼预测、风险评估和消息发布五个阶段。对每个阶段记录输入时间戳、输出时间戳及处理频率，便可估计其均值时延和波动范围。对于风险评估部分，还应关注它与视觉投影链路之间是否形成等待关系，即系统是否需要同时等待吊装构件状态和动态目标预测结果才能完成一次完整预警更新。若确实存在这种同步依赖，则动态链路的抖动会进一步放大端到端预警延迟。

由于该部分当前仍缺少与视觉模块同等级的内置计时统计，因此正文应写为“在现有预测模块基础上补充实验统计”。这种写法并不回避系统现状，而是说明论文已经识别出完整实时性分析所需的关键环节。后续数据可统一记为 \textbf{【待补充：动态检测、聚类、预测与风险评估阶段耗时】}。

\subsection{端到端时延预算与工程实时性分析}\label{subsec_ch5_end_to_end_latency}

单个模块的时延统计只有在被纳入统一预算后，才能真正回答“系统是否满足吊装实时预警需求”这一核心问题。因此，本小节需要从传感器采样开始，综合视觉推理、点云投影、动态候选生成、轨迹预测、风险判级和控制建议输出，建立完整的端到端时延预算框架。设总时延为 $T_{\mathrm{all}}$，则可写为
\begin{equation}
T_{\mathrm{all}}=T_{\mathrm{sensor}}+T_{\mathrm{vision}}+T_{\mathrm{projection}}+T_{\mathrm{dynamic}}+T_{\mathrm{risk}}+T_{\mathrm{publish}}.
\end{equation}
这一表达虽然简化，但能够清楚说明系统延迟并不是由单一模型决定，而是由多阶段处理共同累积形成。

在工程评估中，端到端实时性的判断标准应与吊装作业的风险演化速度相关，而不是只看系统能否达到某个固定帧率。若目标接近速度较高，则系统需要更小的总时延才能保留足够的干预时间；若场景主要是慢速摆动或缓慢靠近，则时延要求可以适当放宽。因此，本小节还应结合动态场景实验中的接近速度范围，对时延预算进行场景化解释。对于当前阶段，可先给出预算框架与评估思路，并将具体时延表和安全余量分析标记为 \textbf{【待补充：端到端时延预算表与实时性结论】}。

\begin{figure}[htbp]
	\centering
	\fbox{\parbox[c][5.2cm][c]{0.88\textwidth}{\centering
	占位图：端到端时延预算分解图\\[0.5em]
	建议内容：传感器采样、视觉推理、点云投影、动态检测、风险评估、消息发布各阶段耗时的堆叠条形图或流程预算图，并标出总时延与安全余量。}}
	\caption{端到端时延预算与实时性分析占位图}
	\label{fig:ch5_latency_budget_placeholder}
\end{figure}

\section{系统稳定性与工程应用可行性分析}\label{sec_ch5_stability_feasibility}

\subsection{多模块协同运行稳定性分析}\label{subsec_ch5_stability}

吊装安全预警系统的实际可用性，除了取决于单项算法精度外，更取决于多模块在长时间运行中的协同稳定性。对于本文系统而言，稳定性问题主要体现在四个方面。第一，图像与点云之间的时间对齐是否持续稳定，若同步窗口频繁失配，则投影结果和空间定位会直接波动。第二，TF 变换链是否连续可靠，若某一环节坐标变换失效，则投影链路将无法正常建立。第三，边缘端计算资源是否稳定，若推理时延和消息处理时延出现明显波动，则预警更新周期会失去可预测性。第四，动态点检测和聚类结果是否足够稳定，若动态候选状态频繁闪烁，则趋势预测和风险升级规则也会随之失去稳定基础。

因此，本小节的分析重点应从“单帧效果”转向“连续运行表现”。例如，可通过长时间 rosbag 回放或在线连续运行，统计消息掉帧率、同步失败次数、TF 查询失败次数、分割结果中断次数以及动态轨迹 ID 频繁切换的情况。通过这些指标，系统的稳定性不再停留在定性描述，而能够形成较为明确的工程评估依据。对于尚未整理完成的长时间运行数据，可保留为 \textbf{【待补充：连续运行日志统计与异常情况分析】}。

\subsection{工程部署可行性与当前局限}\label{subsec_ch5_feasibility_limit}

从工程部署角度看，本文方案具有若干明确优势。首先，系统以模块化 ROS 节点组织视觉、点云、动态检测和预测链路，便于单模块替换和分阶段调试。其次，YOLO26 已完成 Jetson 边缘端部署，说明视觉感知部分具备实际嵌入式运行条件。再次，点云投影回查与动态预测链路均已具备实现基础，使系统不仅能“看见目标”，还能够构建三维安全语义。上述特点共同说明，本研究路线具备较好的工程延展性。

与此同时，系统也存在明确局限。其一，训练参数、分割精度和部分实验统计尚未完全整理，导致部分章节仍需通过 \textbf{【待补充】} 占位。其二，静态距离分层预警、TTC 风险升级和控制接口闭环尚未在仓库中形成完整实现，因此当前系统仍更接近“感知与预警原型”，而非实机控制级成品。其三，复杂工况下的参数敏感性，例如同步窗口、下采样比例、阈值设定和预测时域等，仍需通过更多实验进行系统标定。明确这些局限并不削弱论文贡献，反而能够更真实地说明本文已经完成的工作边界以及后续可继续推进的方向。